\Needspace{9\baselineskip}\begin{center}\textit{Solving the congruence \(xx\equiv A\) by the method of exclusion.}\end{center}


\Needspace{6\baselineskip}
\Article{319}
We treated the \textit{solvability} of the congruence \(xx\equiv A\; (\mathrm{mod}.\ m)\), which is equivalent to the indeterminate equation \(xx=A+my\), in Section IV (Article \(146\)) in such a way that nothing further seems to be desired. As for finding the unknown itself, however, we have already observed above (Article \(152\)) that indirect methods are far preferable to direct ones. If \(m\) is a prime number (to which case the others are easily reduced), we could use the table of indices I for this purpose, combined with III according to the observation in Article \(316\), as we showed more generally in Article \(60\); but this method would be restricted to the limits of the table. For these reasons, we hope that lovers of arithmetic will welcome the following general and convenient method.

First of all, we observe that it is enough to find only those values of \(x\) that are positive and no greater than \(\frac12m\), since every other value is congruent modulo \(m\) to one of these or its negative. For such a value of \(x\), the value of \(y\) must lie between the bounds \(-\frac Am\) and \(\frac14m-\frac Am\). The method that immediately suggests itself would therefore consist in taking each value of \(y\) within these bounds (denoting their collection by \(\Omega\)), calculating the value of \(A+my\), which we shall denote by \(V\), and retaining only those for which \(V\) is a square. When \(m\) is a small number (\textit{e.g.} below \(40\)), this trial is so brief that there is hardly any need to shorten it; but when \(m\) is large, the work can be shortened as much as we please by the following \textit{method of exclusion}.

\Needspace{6\baselineskip}
\Article{320}
Let \(E\) be any integer coprime to \(m\) and greater than \(2\), and let all its distinct quadratic nonresidues (\textit{i.e.} incongruent modulo \(E\)) be \(a,\allowbreak b,\allowbreak c\), etc.;
\SourceLeaf{394}
finally, let the roots of the congruences
\[A+my\equiv a,\quad A+my\equiv b,\quad A+my\equiv c\quad\text{etc. modulo }E\]
be \(\alpha,\allowbreak \beta,\allowbreak \gamma\), etc., which may all be taken positive and less than \(E\). If, then, \(y\) is assigned a value congruent to one of these numbers \(\alpha,\allowbreak \beta,\allowbreak \gamma\), etc., modulo \(E\), the resulting value of \(V=A+my\) will be congruent to one of \(a,\allowbreak b,\allowbreak c\), etc., and hence a nonresidue of \(E\); it therefore cannot be a square. Thus, from \(\Omega\), all the numbers of the forms \(Et+\alpha,\allowbreak Et+\beta,\allowbreak Et+\gamma\), etc., can at once be excluded as useless, and it will be enough to try the remaining numbers, whose collection we shall call \(\Omega'\). In this operation, the number \(E\) may be called an \textit{exclusion number}.

Now, taking another suitable number \(E'\) as an exclusion number, we shall find in exactly the same way as many numbers \(\alpha',\allowbreak \beta',\allowbreak \gamma'\), etc., as it has distinct quadratic nonresidues, to which \(y\) cannot be congruent modulo \(E'\). We may therefore again remove from \(\Omega'\) all the numbers of the forms \(E't+\alpha',\allowbreak E't+\beta',\allowbreak E't+\gamma'\), etc. We can continue in this way, using further exclusion numbers, until the collection of numbers in \(\Omega\) has been reduced so much that actually trying all those left seems no more difficult than making further exclusions.

\textit{Example.} Given the equation \(xx=22+97y\), the bounds for the values of \(y\) will be \(-\frac{22}{97}\) and \(24\frac14-\frac{22}{97}\), so that, since the value \(0\) is obviously useless, \(\Omega\) will contain the numbers \(1,\allowbreak 2,\allowbreak 3\ldots24\). For \(E=3\) there is just one nonresidue, \(a=2\), giving \(\alpha=1\); hence we must exclude from \(\Omega\) all numbers of the form \(3t+1\), and the number of those remaining, \(\Omega'\), will be \(16\). Similarly, for \(E=4\) we have \(a=2,\allowbreak b=3\), giving \(\alpha=0,\allowbreak \beta=1\); hence the numbers of the forms \(4t\) and \(4t+1\) must be rejected, leaving these eight: \(2,\allowbreak 3,\allowbreak 6,\allowbreak 11,\allowbreak 14,\allowbreak 15,\allowbreak 18,\allowbreak 23\). Likewise, for \(E=5\), we find that the numbers of the forms \(5t\) and \(5t+3\) must be rejected, leaving \(2,\allowbreak 6,\allowbreak 11,\allowbreak 14\). The exclusion number \(6\) would remove numbers of the forms \(6t+1\) and \(6t+4\), but these, which coincide with the numbers of the form \(3t+1\), are already absent. The exclusion number \(7\) removes the numbers of the forms \(7t+2,\allowbreak 7t+3,\allowbreak 7t+5\), leaving \(6,\allowbreak 11,\allowbreak 14\). Substituting these for \(y\) gives, respectively, \(V=604,\allowbreak 1089,\allowbreak 1380\), of which only the second value is a square; hence \(x=\mp33\).
\SourceLeaf{395}

\Needspace{6\baselineskip}
\Article{321}
The operation with the exclusion number \(E\) removes, from the values of \(V\) corresponding to the values of \(y\) in \(\Omega\), all those that are quadratic nonresidues of \(E\), while leaving the residues of the same number untouched. It is therefore easy to see that using the exclusion numbers \(E\) and \(2E\) makes no difference if \(E\) is odd, since in this case \(E\) and \(2E\) have the same residues and nonresidues. Hence, if the numbers \(3,\allowbreak 4,\allowbreak 5\), etc., are used successively as exclusion numbers, the numbers \(6,\allowbreak 10,\allowbreak 14\), etc., which are twice odd numbers, must be skipped as superfluous. It is also clear that the two operations with the exclusion numbers \(E,\allowbreak E'\) remove all those values of \(V\) that are nonresidues of both \(E,\allowbreak E'\) or of one of them, leaving those that are residues of both. Now, when \(E\) and \(E'\) have no common divisor, all the numbers removed are nonresidues of the product \(EE'\), and all those left are its residues. It is therefore evident that using the exclusion number \(EE'\) has exactly the same effect in this case as using both \(E,\allowbreak E'\), and is thus superfluous after these have been used. We may therefore also skip all exclusion numbers that can be resolved into two coprime factors, and it will be enough to use prime numbers that do not divide \(m\), or powers of primes. Finally, it is clear that, after the exclusion number \(p^\mu\), a power of the prime \(p\), has been used, the exclusion number \(p\) or \(p^\nu\), where \(\nu<\mu\), becomes superfluous. For, since \(p^\mu\) has left only its own residues among the values of \(V\), certainly no nonresidues of \(p\) or of any lower power \(p^\nu\) can remain. If, however, \(p\) or \(p^\nu\) has already been used before \(p^\mu\), the latter can evidently remove only those values of \(V\) that are at once residues of \(p\) (or \(p^\nu\)) and nonresidues of \(p^\mu\); hence it will be enough to take only nonresidues of \(p^\mu\) of this kind for \(a,\allowbreak b,\allowbreak c\), etc.

\Needspace{6\baselineskip}
\Article{322}
The calculation of the numbers \(\alpha,\allowbreak \beta,\allowbreak \gamma\), etc., corresponding to any given exclusion number \(E\) is greatly shortened by the following observations. Let \(\mathfrak A,\allowbreak \mathfrak B,\allowbreak \mathfrak C\), etc., be roots of the congruences \(my\equiv a,\allowbreak my\equiv b,\allowbreak my\equiv c\), etc., \((\mathrm{mod.}\ E)\), and \(k\) a root of \(my\equiv-A\); it will then be clear that \(\alpha\equiv\mathfrak A+k,\allowbreak \beta\equiv\mathfrak B+k,\allowbreak \gamma\equiv\mathfrak C+k\), etc. Now, if we really had to obtain \(\mathfrak A,\allowbreak \mathfrak B,\allowbreak \mathfrak C\), etc., by solving these congruences, this way of finding \(\alpha,\allowbreak \beta,\allowbreak \gamma\), etc., would certainly be no shorter than the one shown above; but this is by no means necessary. For, first, if \(E\) is a prime number and \(m\) a quadratic residue of \(E\), it follows from Article \(98\) that \(\mathfrak A,\allowbreak \mathfrak B,\allowbreak \mathfrak C\), etc.,
\SourceLeaf{396}
which are the values of the expressions \(\frac am,\allowbreak \frac bm,\allowbreak \frac cm\), etc., \((\mathrm{mod.}\ E)\), are distinct nonresidues of \(E\), and therefore coincide exactly with \(\alpha,\allowbreak \beta,\allowbreak \gamma\), etc., apart from their order, which does not matter here. If, under the same assumption, \(m\) is a nonresidue of \(E\), the numbers \(\mathfrak A,\allowbreak \mathfrak B,\allowbreak \mathfrak C\), etc., will coincide with all the quadratic residues except \(0\). If \(E\) is the square of an odd prime number, \(=pp\), and \(p\) has already been used as an exclusion number, it is enough, by the preceding article, to take for \(a,\allowbreak b,\allowbreak c\), etc., those nonresidues of \(pp\) that are residues of \(p\), \textit{i.e.} the numbers \(p,\allowbreak 2p,\allowbreak 3p\ldots pp-p\) (namely, all numbers below \(pp\) except \(0\) that are divisible by \(p\)). It is then easy to see that \(\mathfrak A,\allowbreak \mathfrak B,\allowbreak \mathfrak C\), etc., must give exactly the same numbers, only in a different order. Similarly, if, after the exclusion numbers \(p\) and \(pp\) have been used, we put \(E=p^3\), it will be enough to take for \(a,\allowbreak b,\allowbreak c\), etc., the products of the individual nonresidues of \(p\) by \(pp\); then \(\mathfrak A,\allowbreak \mathfrak B,\allowbreak \mathfrak C\), etc., will give either the same numbers or the products of \(pp\) by the individual residues of \(p\) except \(0\), according as \(m\) is a residue or a nonresidue of \(p\). In general, taking for \(E\) any power of a prime number, say \(p^\mu\), after all lower powers have been used, we shall obtain for \(\mathfrak A,\allowbreak \mathfrak B,\allowbreak \mathfrak C\), etc., the products of \(p^{\mu-1}\) either by all numbers less than \(p\), always excepting \(0\), when \(\mu\) is even; or by all nonresidues of \(p\) less than \(p\), when \(\mu\) is odd and \(mRp\); or by all residues, when \(mNp\). — If \(E=4\), and hence \(a=2,\allowbreak b=3\), we have for \(\mathfrak A,\allowbreak \mathfrak B\) either \(2\) and \(3\) or \(2\) and \(1\), according as \(m\equiv1\) or \(\equiv3\; (\mathrm{mod}.\ 4)\). If, after using the exclusion number \(4\), we put \(E=8\), we have \(\alpha=5\), giving \(\mathfrak A\) the values \(5,\allowbreak 7,\allowbreak 1,\allowbreak 3\), according as \(m\equiv1,\allowbreak 3,\allowbreak 5,\allowbreak 7\; (\mathrm{mod}.\ 8)\). In general, however, if \(E\) is any higher power of two, say \(2^\mu\), after the lower powers have been used, we must put \(a=2^{\mu-1},\allowbreak b=3.2^{\mu-2}\) when \(\mu\) is even, giving \(\mathfrak A=2^{\mu-1},\allowbreak \mathfrak B=3.2^{\mu-2}\) or \(=2^{\mu-2}\), according as \(m\equiv1\) or \(\equiv3\); but when \(\mu\) is odd, we must put \(a=5.2^{\mu-3}\), giving \(\mathfrak A\) equal to the product of \(2^{\mu-3}\) by \(5,\allowbreak 7,\allowbreak 1\) or \(3\), according as \(m\equiv1,\allowbreak 3,\allowbreak 5\) or \(7\; (\mathrm{mod}.\ 8)\).

Those familiar with the subject will, moreover, readily devise an apparatus by which the useless values of \(y\) can be removed \textit{mechanically} from \(\Omega\), once the numbers \(\alpha,\allowbreak \beta,\allowbreak \gamma\), etc., have been calculated for as many exclusion numbers as seem necessary; but we cannot discuss this here, any more than other devices for shortening the work.
